% TOP_PROBLEM: {"id": 30007148, "title": "Metric classification of Blaschke manifolds", "category_id": 6, "category": {"id": 6, "name": "geometry", "display_name": "Geometry", "description": "Euclidean and non-Euclidean geometry, geometric structures.", "slug": "geometry", "order_index": 6, "created_at": "2026-07-31T15:26:25.671Z"}, "status": "open"}
% FIRST_POSTED: 2026-10-06
% ATTEMPT_STATUS: solved
% OPENAI_CLAIM: {"schema_version":1,"source_repo":"https://github.com/openai/math","source_commit":"adc7f1241b42e322a6451854ab7e4b4c146bf78a","release_date":"2026-10-06","match":"full","resolution":"proved","independently_reviewed":false,"second_pass":"agreed","families":[{"family":"344","title":"The metric Blaschke conjecture","subject":"Differential geometry","manuscripts":[{"title":"The metric Blaschke theorem","pdf_path":"preprints/The-metric-Blaschke-theorem-September-23-2026/paper.pdf","readme_path":"preprints/The-metric-Blaschke-theorem-September-23-2026/README.md","date":"2026-09-23","theorem_ref":"preprints/The-metric-Blaschke-theorem-September-23-2026/build/sections/00-introduction.tex#thm:blaschke"}],"lean":null,"reasoning_trace":null}]}
\documentclass[11pt]{article}
\usepackage[margin=25mm]{geometry}
\usepackage{fontspec}
\setmainfont{Latin Modern Roman}
\usepackage{amsmath,amssymb,mathtools}
\usepackage{unicode-math}
\setmathfont{Latin Modern Math}
\usepackage{enumitem}
\usepackage{xurl,hyperref,bookmark}
\hypersetup{colorlinks=true,urlcolor=blue}
\setcounter{secnumdepth}{0}
\setlength{\parindent}{0pt}
\setlength{\parskip}{5pt}
\setlength{\emergencystretch}{3em}
\begin{document}
\section*{Metric classification of Blaschke manifolds}

\subsection{OpenAI's claim}
\textit{Metric Blaschke rigidity}: The manuscript claims that a closed connected manifold with injectivity radius equal to diameter is a scaled standard compact rank-one symmetric space. The assertion includes all positive dimensions and model types.

\subsection{Outline of the argument}
\textit{Metric Blaschke rigidity}: Normal Jacobi fields along a closed geodesic are encoded as Lagrangian planes. A fixed graph chart factors the endpoint determinant using a nonsingular block and its Schur complement, with the latter controlled by the cut-time intersection data. Sharp volume comparison and equality analysis combine with the corresponding upper bound, forcing the metric to equal its model.

\subsection{Scope relative to this problem}
\textit{Metric Blaschke rigidity}: The manuscript claims that a closed connected manifold with injectivity radius equal to diameter is a scaled standard compact rank-one symmetric space. The assertion includes all positive dimensions and model types. The parameters, regularity assumptions and normalization in this question are retained in applying that assertion.

With the hypotheses and conventions above, the claimed result addresses the whole mathematical assertion on this page.

\subsection{Formal verification}
\textit{Metric Blaschke rigidity}: The pinned catalogue lists no Lean formalization for this result.

This site has not run the Lean development. These links describe OpenAI’s selected formal statements and should be read with their stated scope.

\subsection{Status caveat}
OpenAI reports that an internal model produced the manuscripts. Their arguments have not been independently verified as proofs by this site. Any portion outside the documented formal scope may contain mistakes, and the catalogue label reports the announced mathematical claim.

\subsection{Sources}
\begin{itemize}
\item OpenAI, \emph{The metric Blaschke theorem}, 2026-09-23; citation key \texttt{OAI:The-metric-Blaschke-theorem-September-23-2026}. \href{https://github.com/openai/math/blob/adc7f1241b42e322a6451854ab7e4b4c146bf78a/preprints/The-metric-Blaschke-theorem-September-23-2026/paper.pdf}{PDF (GitHub)}; \href{https://github.com/openai/math/raw/adc7f1241b42e322a6451854ab7e4b4c146bf78a/preprints/The-metric-Blaschke-theorem-September-23-2026/paper.pdf}{PDF (direct)}; \href{https://github.com/openai/math/blob/adc7f1241b42e322a6451854ab7e4b4c146bf78a/preprints/The-metric-Blaschke-theorem-September-23-2026/README.md}{Manuscript page and citation}; \href{https://github.com/openai/math/blob/main/preprints/The-metric-Blaschke-theorem-September-23-2026/paper.pdf}{Latest version}.
\item \href{https://github.com/openai/math/blob/adc7f1241b42e322a6451854ab7e4b4c146bf78a/CONTENTS.md}{OpenAI catalogue entry}.
\end{itemize}
\end{document}
